\section{Compilation}\label{ch5sec:compilation}

The template can be compiled with \texttt{pdflatex}, \texttt{lualatex}, or \texttt{xelatex}. The default is \texttt{pdflatex} because it is supported by arXiv and works in restricted environments. Compilation is orchestrated by \texttt{latexmk} through either the Makefile or the \texttt{.latexmkrc} file.

\paragraph{Using the Makefile.} Running \texttt{make} builds \texttt{build/main.pdf}. Running \texttt{make clean} removes auxiliary files, and \texttt{make cleanall} removes the entire build directory.

\begin{table}[h]
\caption{Supported compilation engines}\label{ch5tab:demo-datasets}
\centering
\begin{tabular}{l l l}
\toprule
\thead{Engine} & \thead{Best for} & \thead{Notes} \\
\midrule
\texttt{pdflatex} & arXiv compatibility & Default engine \\
\texttt{lualatex} & Modern OpenType fonts & Flexible font handling \\
\texttt{xelatex} & System fonts & Similar to \texttt{lualatex} \\
\bottomrule
\end{tabular}
\end{table}

\paragraph{Using latexmk directly.} You can also run \texttt{latexmk main.tex}. The \texttt{.latexmkrc} file configures \texttt{pdflatex}, enables shell escape for \texttt{minted}, and redirects output to \texttt{build/}. Table~\ref{ch5tab:demo-datasets} compares the available engines.
